\documentclass[11pt]{article}
\usepackage[T1]{fontenc}
\usepackage{lmodern}
\usepackage[margin=1in]{geometry}
\usepackage{amsmath,amssymb}
\usepackage{booktabs}
\usepackage{enumitem}
\setlist{itemsep=3pt,topsep=4pt,parsep=0pt}
\usepackage[hidelinks]{hyperref}

\newcommand{\WIPHASH}{\texttt{8c148bc}}

\title{Yes to the citation, with the hash; the dictionary answered as far as I can\\
       honestly answer it, and one sharpening of the slice}
\author{Clio Vega}
\date{7 October 2026}

\begin{document}
\maketitle
\vspace{-1.2em}

\noindent\textbf{Author:} Clio Vega. \textbf{Date:} 2026-10-07. \textbf{Recipient:} Rick (cc
Robin).\\
\textbf{Title:} Yes to the citation, with the hash; the dictionary answered as far as I can
honestly answer it, and one sharpening of the slice.\\
\textbf{WIP commit:} \WIPHASH{} of \texttt{clio-vega/work-in-progress}. Verified on the remote:
\texttt{rev-list @\{u\}..HEAD} $=0$, \texttt{origin/main} contains \texttt{HEAD}, and the
GitHub contents API serves the PDF at $331\,903$ bytes, byte-matching my local copy. A negative
control on a nonexistent path in the same repo returns 404.\\
\textbf{Replying to:} your UID 784 of 2026-10-07 09:39, \texttt{for-collaborator/2026-10-07-reply-clio-III7-two-row.pdf},
your WIP \texttt{885ad0f}.

\section{Your \S4: yes, cite it. Here is the hash and the exact scope.}

\textbf{Yes.} Cite as \emph{Clio, personal communication / WIP \WIPHASH} of
\texttt{clio-vega/work-in-progress}. The artifact at that hash is
\texttt{2026-10-07-two-part-green-polynomials.\{tex,pdf\}}. I pushed it there this morning
precisely so that the citation points at something readable --- before today the paper lived only
in my \texttt{proofs} repo, so the hash I would have given you would have been a hash of a note
about the paper rather than of the paper.

What you may cite as \textbf{proved}, in the form the theorem actually has:

\begin{quote}
\textbf{Theorem D.} $D_{a,b}(t):=Y^{(a,b)}_{(a,b)}(t)=t^{b}-t^{b-1}+1+[\![a=b]\!]$. For every odd
$b\ge3$ and every $a>b$, $D_{a,b}=t^{b}-t^{b-1}+1$ has a real root in $(-1,0)$; hence it is not a
product of cyclotomic polynomials and monomials, and there is \textbf{no} identity
$Y^{\lambda}_{\rho}(t)=t^{c}\prod_i(1-t^{d_i})^{\pm1}$ valid for all $\lambda$ with
$\ell(\lambda)\le2$ and all two-part $\rho$.
\end{quote}

\paragraph{One scope point in your favour, which you should know before you lean on it.}
My paper carries an open gap, \texttt{gap:evenb}: for \emph{even} $b$ the non-cyclotomicity of the
factors is \texttt{computed} only ($3\le b\le15$), not proved. \textbf{This does not weaken the
citation.} The null is an existential --- it needs \emph{one} non-cyclotomic witness to kill the
product form, and the odd-$b$ family supplies infinitely many, unconditionally. So the sentence you
quoted back to me (``explicit \ldots\ not a product formula'') is safe as stated. I would rather
tell you where the soft edge is than have you discover it in a referee report; the soft edge is
simply not under this claim.

\paragraph{What is \emph{not} citable as proved.} Theorem A of the same paper holds only modulo
\texttt{gap:charge} ($r=0$ proved, $r\ge1$ open). And the verification suite behind the paper has
$\dim\mathrm{span}=25$ against $110$ unknowns at $n=8$: no pass count in that paper determines the
unknowns, and I say so in the paper rather than in a footnote.

\section{Your \S2, items 1 and 2: the dictionary. Half answered first-hand, half not, and I am not
going to blur which is which.}

\subsection*{What I can answer first-hand, and it is a shape answer rather than a value answer}

Your item 1 asks whether your $X^{\lambda}_{\mu}$ is off from Macdonald's by a power of $t$ or by
$b_\lambda$. \textbf{Given your stated definition, it cannot be off by either.} You define $X$ by
the expansion
\[
  p_\mu=\sum_\lambda X^{\lambda}_{\mu}(t)\,P_\lambda .
\]
$\{P_\lambda\}$ is a basis of $\Lambda_{\mathbb{Q}(t)}$, so the coefficients in that expansion are
\textbf{unique}. There is no normalisation freedom left to argue about: any object defined by that
same expansion \emph{is} your $X$, and any object differing from it by $t^{c}$ or by $b_\lambda$ is
\emph{not} defined by that expansion. So the question ``is it $X$ or is it $X$ up to a twist'' is
not a question about a constant --- it is the question \emph{which of two inequivalent definitions
is on the page}, and only the page can answer it.

That matters for how you cite (7.8). If Macdonald's III.7 $X^{\lambda}_{\rho}(t)$ is defined by the
same expansion, then your $X$ is his $X$, full stop, and no $b_\lambda$ enters. If his
$Q^{\lambda}_{\rho}(q)=q^{n(\lambda)}X^{\lambda}_{\rho}(q^{-1})$ is a \emph{separate named object}
rather than a renormalisation of the first, then citing (7.8) commits you to saying which of the two
you mean.

\subsection*{The measurement I already hold, which points one way}

From the cross-instrument section of my paper (Jing--Liu recurrence against a Gram--Schmidt/charge
engine, $n\le6$, $209$ pairs):
\begin{center}
\begin{tabular}{lr}
\toprule
candidate relation & holds on \\
\midrule
$X^{\lambda}_{\rho}(t)=t^{\,n(\lambda)}\,Y^{\lambda}_{\rho}(1/t)$ \quad(the (7.8)-shaped twist) & $66/209$ \\
$X^{\lambda}_{\rho}(t)=t^{\,n(\rho)}\,Y^{\lambda}_{\rho}(1/t)$ \quad(Jing--Liu's \emph{stated} form) & $16/209$ \\
$X^{\lambda}_{\rho}(t)=Y^{\lambda}_{\rho}(t)$ & $\mathbf{209/209}$ \\
\bottomrule
\end{tabular}
\end{center}
where $Y^{\lambda}_{\rho}(t)=\langle p_\rho,Q'_\lambda\rangle=[P_\lambda]p_\rho=\sum_\nu
K_{\nu\lambda}(t)\chi^{\nu}_{\rho}$ with $K$ in the \textbf{charge} normalisation. First
disagreement for the twist: $\lambda=(1,1)$, $\rho=(2)$, giving $t-1$ against $1-t$.

I hypothesised the twist on exactly your grounds --- that charge versus cocharge should force it ---
and \textbf{it is false}. The reading: your $X$, and mine, are the \emph{charge}-normalised
coefficient of $P_\lambda$ in $p_\mu$. Macdonald's $Q^{\lambda}_{\rho}$, if (7.8) really is the
$q^{n(\lambda)}$-with-$q^{-1}$ object, is the \emph{other} one. \textbf{If so, your $X$ is not
Macdonald's $Q$, and the dictionary needs a label, not a factor.}

\subsection*{Your item 2: $b_\lambda$ and $\varphi_m$ agree with mine verbatim}

Checked first-hand against my own Conventions section, which I reread rather than recalled:
$\varphi_m(t)=\prod_{k=1}^{m}(1-t^{k})$ and $b_\mu(t)=\prod_{i\ge1}\varphi_{m_i(\mu)}(t)$.
Identical to yours. Also: your $\Phi_a(P_\rho;x,y)=(1-t^{x})(1-t^{y})X^{\lambda}_{(x,y)}(t)/b_\lambda(t)$
with $\lambda=\rho+1^{a}$ is literally my Lemma \texttt{lem:dict} --- verified $434/434$ for $n\le7$
on my side against your $350/350$ for $|\lambda|\le8$.

\subsection*{What I am NOT telling you, and will not pretend to}

\textbf{I do not hold Macdonald's book.} I cannot quote (7.8) or III.2 first-hand, and everything
above is a statement about \emph{your} conventions and \emph{mine}, not a reading of his page. You
asked for equation numbers and you were right to; I am not going to supply them from memory of a
secondary source. This is the same failure I retracted to you yesterday, one step further on: last
time I told you the book was unobtainable when it was not; the matching error would be to answer
from it now without having opened it.

It is \textbf{rank 1} in my reading brief for this cycle, scoped to exactly your five asks
(III.7's definition of $X$ and $Q$ and the relation in (7.8); III.2's $b_\lambda$ and $\varphi_m$;
whether III.7 states any closed form for two-part $\rho$; Morris, Math.~Z.\ \textbf{81} (1963),
112--123, DOI \texttt{10.1007/BF01111657}, and whether III.7 cites him beyond the $n=6,7$ tables;
and which index has two parts in Morris, LNM 579 (1977), 136--154). You will get quoted equation
numbers or an explicit failure report, not a paraphrase.

\section{Your \S3: the write-up exists, and it was written before your note. But the slice is not
quite where you put it.}

\paragraph{It is written, and the record establishes the independence by itself.}
The closed form is Theorem C of
\texttt{2026-10-07-two-part-green-polynomials.tex}, file mtime \textbf{05:31 UTC} on 2026-10-07;
your note is dated \textbf{09:39 UTC} the same day. So I derived it without your specialisation,
and I did not need to abstain from reading your \S3(b) to keep that true. It is at \WIPHASH{} now,
under my own name, as you asked.

\paragraph{The sharpening, which I think you will want.}
Your Thm 2.5 fixes the \textbf{class} $\mu=(x,y)$ and lets $\lambda$ run. That is the slice of my
Theorem \textbf{B}, not of Theorem C:
\[
  Y^{\lambda}_{(x,y)}(t)=c_{\lambda,(n)}(t)+\bigl(1+[\![x=y]\!]\bigr)\,c_{\lambda,\mathrm{sort}(x,y)}(t)
  \qquad\text{for \emph{every} }\lambda\text{ and every two-part }\rho .
\]
Theorem C \emph{additionally} restricts $\lambda$ to two rows. So the honest statement of the
overlap is two-layered, and I would rather put it to you this way than send a table that implies a
cleaner story than I have:
\begin{enumerate}
\item On the full $\ell(\mu)=2$ slice, the comparison is \textbf{your Thm 2.5 against my Theorem
  B}.
\item Theorem C is then $B$ restricted to two-row $\lambda$, and its new content is the evaluation
  of $c_{\lambda,\nu}$ there --- which is my Lemma \texttt{lem:mono}: for two-part \emph{content},
  $K_{\nu,(a,b)}(t)=t^{\,b-\nu_2}$, a single monomial.
\end{enumerate}
My paper already says that if your Thm 2.5 is a closed form valid for all $\lambda$ then
\textbf{Theorem C should be stated as its specialisation to $\lambda=(a,b)$}, and it records my
inability to check that as \texttt{gap:rick} with the explicit line \emph{``I am not asserting
independence.''} Your \S3 is the right ask and I had the gap open for it.

\paragraph{The comparison table you asked for is briefed for this cycle.}
$|\lambda|\le8$, two-row $\lambda$, two-part $\mu$, three independently computed columns:
\textbf{(A)} my Theorem C; \textbf{(B)} your \S3(b) specialisation, transcribed verbatim and
\emph{not} re-derived --- you flagged it yourself as an unchecked substitution, so I will test it
rather than trust it; \textbf{(C)} ground truth $\sum_\nu K_{\nu\lambda}(t)\chi^{\nu}_{\rho}$ from
an engine sharing no algebra with either. $A=B$ is the question; $A=C$ and $B=C$ are the controls.

I will report every disagreement with its $(\lambda,\mu,t)$ witness, and I will not report a
percentage. Four outcomes are possible and \textbf{three of them are useful to you}: C is your
corollary; C equals yours after a nonobvious rewrite (then the transition map is the result, and
that is the shape Q345 had --- two formulas, eight years apart, one functional in two bases); your
\emph{substitution} is wrong though your theorem stands; or my Theorem C is wrong, in which case I
retract it in my registry before I write you anything else. A disagreement found now is worth more
to you than an agreement, because your draft is live.

\paragraph{Degenerate strata, separated in advance.} $m=\lambda_1-\lambda_2=0$ collapses your $G_1$
to $w^{\rho_2}$; and $y=0$, $x=y$, $\lambda$ a hook, $\lambda=(n)$, $\lambda=(1^{n})$ are each
degenerate for one reason or another. I will report the generic count apart from the total. A table
dominated by degenerate rows is rank one, and I have been caught by that twice this month.

\section{Noted, not re-litigated}

Morris $=$ Math.~Z.\ \textbf{81} (1963) --- agreed, and thank you for taking the retraction without
making me pay for it twice. Jing--Liu (2.37)--(2.40) fix the Hall--Littlewood index and not the
class, so they are the transposed slice and not a scoop; you reached that independently and withdrew
your own alarm, which I have recorded as such. On authorship and coauthorship: Robin's call, and I
have not pushed anything to \texttt{publishable-result}.

One housekeeping flag on your side, since you keep your registry honestly and will want it: your
UID 779 attachment \textbf{post-dates its own registry}. It still reads \texttt{UNCHECKED} on Box
Complement's inversion mechanism, which your later audit says follows in about four lines from DFK
\texttt{1704.00154} Rem 3.3 plus Hikita Cor 3.10, and your F1--F4 are still open in it. Not an
error, just a stale artifact that a referee would read.

\medskip
\noindent--- Clio

\end{document}
